\begin{figure}[htbp]
	\centering
	\begin{tikzpicture}[>=Stealth, line width=0.8pt, font=\small]
		
		\def\ang{25}      % angolo alpha^e in gradi
		\def\Lb{4.2}      % lunghezza trave
		\def\sep{6.5}     % offset pannello destro
		\def\ds{0.72}     % lunghezza frecce spostamento
		\def\fn{0.90}     % lunghezza frecce forza
		\def\fib{0.18}    % distanza fibra inferiore dall'asse
		\def\rarc{0.32}   % raggio arco rotazione (cinematica)
		\def\rmarc{0.42}  % raggio arco momento (statica)
		
		%%=============================================================
		%% PANNELLO (a) — CINEMATICA
		%%=============================================================
		
		\coordinate (Pi) at (0, 0);
		\coordinate (Pj) at ({\Lb*cos(\ang)}, {\Lb*sin(\ang)});
		
		%% Fibra inferiore (grigia tratteggiata)
		\coordinate (Pif) at ({\fib*sin(\ang)},               {-\fib*cos(\ang)});
		\coordinate (Pjf) at ({\Lb*cos(\ang)+\fib*sin(\ang)}, {\Lb*sin(\ang)-\fib*cos(\ang)});
		\draw[dashed, gray!70, line width=1.0pt] (Pif) -- (Pjf);
		
		%% Assi globali
		\draw[->] (-1.0,-1.2) -- ++(1.3,0) node[right]{$x$};
		\draw[->] (-1.0,-1.2) -- ++(0,1.2) node[above]{$y$};
		
		%% Segmento orizzontale grigio (riferimento per alpha^e)
		\draw[gray!60, thin] (Pi) -- ++(2.0,0);
		
		%% Arco alpha^e grigio (raggio ampio, a destra della fibra)
		\draw[->, gray!60] (1.50,0) arc (0:\ang:1.50);
		\node[gray!80] at ({1.82*cos(\ang/2)},{0.18}) {$\alpha^e$};
		
		%% Asta (linea d'asse)
		\draw[line width=1.4pt] (Pi) -- (Pj);
		
		%% Nodi (cerchietti vuoti)
		\filldraw[fill=white, draw=black, line width=1pt] (Pi) circle(3.5pt);
		\filldraw[fill=white, draw=black, line width=1pt] (Pj) circle(3.5pt);
		
		%% Etichette nodi
		\node[below right]  at (Pi) {$P_i$};
		\node[below right] at (Pj) {$P_j$};
		
		%% Spostamenti in P_i
		\draw[->, line width=1.1pt] (Pi) -- ++(\ds,0)  node[below]{$u_i$};
		\draw[->, line width=1.1pt] (Pi) -- ++(0,\ds)  node[right] {$v_i$};
		
		%% Rotazione vartheta_i (arco antiorario, CCW)
		\draw[->] ($(Pi)+(\rarc,0)$) arc (0:270:\rarc);
		\node[above left] at ($(Pi)+(-\rarc+0.02,-0.10)$) {$\vartheta_i$};
		
		%% Spostamenti in P_j
		\draw[->, line width=1.1pt] (Pj) -- ++(\ds,0)  node[below]{$u_j$};
		\draw[->, line width=1.1pt] (Pj) -- ++(0,\ds)  node[right]{$v_j$};
		
		%% Rotazione vartheta_j (arco antiorario, CCW)
		\draw[->] ($(Pj)+(\rarc,0)$) arc (0:270:\rarc);
		\node[above left] at ($(Pj)+(-\rarc+0.02,-0.10)$) {$\vartheta_j$};
		
		%% Etichetta pannello
		\node[font=\small] at ({0.5*\Lb*cos(\ang)},{-1.50}) {$(a)$};
		
		
		%%=============================================================
		%% PANNELLO (b) — STATICA
		%%=============================================================
		
		\coordinate (Pis) at (\sep, 0);
		\coordinate (Pjs) at ({\sep+\Lb*cos(\ang)}, {\Lb*sin(\ang)});
		
		%% Fibra inferiore pannello (b)
		\coordinate (Pisf) at ({\sep+\fib*sin(\ang)},               {-\fib*cos(\ang)});
		\coordinate (Pjsf) at ({\sep+\Lb*cos(\ang)+\fib*sin(\ang)}, {\Lb*sin(\ang)-\fib*cos(\ang)});
		\draw[dashed, gray!70, line width=1.0pt] (Pisf) -- (Pjsf);
		
		%% Assi globali
		\draw[->] ({\sep-1.0},-1.2) -- ++(1.3,0) node[right]{$x$};
		\draw[->] ({\sep-1.0},-1.2) -- ++(0,1.2) node[above]{$y$};
		
		%% Asta
		\draw[line width=1.4pt] (Pis) -- (Pjs);
		
		%% Nodi (cerchietti vuoti)
		\filldraw[fill=white, draw=black, line width=1pt] (Pis) circle(3.5pt);
		\filldraw[fill=white, draw=black, line width=1pt] (Pjs) circle(3.5pt);
		
		%% Etichette nodi
		\node[below right]  at (Pis) {$P_i$};
		\node[below right] at (Pjs) {$P_j$};
		
		%%── Forze in P_j (estremo destro) ──────────────────────────
		%% N_e: direzione +a_e = (cos, sin)
		\draw[->, line width=1.2pt] (Pjs)
		-- ++({cos(\ang)*\fn},{sin(\ang)*\fn})
		node[below right]{$N^e$};
		
		%% T_e: direzione +a_t = (sin, -cos)
		\draw[->, line width=1.2pt] (Pjs)
		-- ++({sin(\ang)*\fn},{-cos(\ang)*\fn})
		node[right]{$T^e$};
		
		%% M_e: arco CCW centrato in P_j
		\draw[->, line width=1.2pt]
		({\sep+\Lb*cos(\ang)+\rmarc*cos(\ang-50)},
		{\Lb*sin(\ang)+\rmarc*sin(\ang-50)})
		arc ({\ang-50}:{\ang+115}:\rmarc)
		node[midway, above]{$M^e$};
		
		%%── Forze in P_i (estremo sinistro, verso opposto) ─────────
		%% N_e: direzione -a_e = (-cos, -sin)
		\draw[->, line width=1.2pt] (Pis)
		-- ++({-cos(\ang)*\fn},{-sin(\ang)*\fn})
		node[below right]{$N^e$};
		
		%% T_e: direzione -a_t = (-sin, +cos)
		\draw[->, line width=1.2pt] (Pis)
		-- ++({-sin(\ang)*\fn},{cos(\ang)*\fn})
		node[left]{$T^e$};
		
		%% (M_e - T_e ell_e): arco CW centrato in P_i (verso opposto a M_e)
		\draw[->, line width=1.2pt]
		({\sep+\rmarc*cos(\ang+115)},
		{\rmarc*sin(\ang+115)})
		arc ({\ang+115}:{\ang-50}:\rmarc)
		node[above, right, font=\footnotesize]
		{$M^e\!-\!T^e\ell^e$};
		
		%% Etichetta pannello
		\node[font=\small] at ({\sep+0.5*\Lb*cos(\ang)},{-1.50}) {$(b)$};
		
	\end{tikzpicture}
	\caption{Trave di telaio $e$ che connette i nodi $P_i$ (iniziale)
		e $P_j$ (finale), inclinata di $\alpha^e$ rispetto all'asse
		globale $x$. La linea tratteggiata grigia indica la fibra
		inferiore e definisce il verso di percorrenza da $P_i$ a $P_j$.
		\emph{(a)}~Cinematica: spostamenti nodali
		$(u_i,\,v_i,\,\vartheta_i)$ e $(u_j,\,v_j,\,\vartheta_j)$
		nel riferimento globale $(x,\,y)$.
		\emph{(b)}~Statica: caratteristiche della sollecitazione
		$N^e$, $T^e$, $M^e$ alle due sezioni di estremità con il
		verso convenzionale positivo; in $P_i$ il momento vale
		$M^e - T^e\ell^e$ in assenza di carichi distribuiti.}
	\label{fig:trave_telaio_cin_stat}
\end{figure}
